\ctufitabstractCZE{Tato studie navazuje na předchozí výzkum, ve kterém byly pro klasifikaci zpravodajských článků nebo jejich shrnutí jako falešných nebo důvěryhodných použity relativně jednoduché klasické modely strojového učení. Cílem této práce je prozkoumat využití pokročilejších architektur — konkrétně modelů založených na transformerech — za účelem zlepšení klasifikační výkonnosti na stejných datových sadách, které byly použity v předchozí studii.}
\ctufitabstractENG{This study extends prior research in which relatively simple, classical machine learning models were employed to classify news articles or summaries as either fake or reliable. The objective of the present work is to investigate the use of more advanced architectures—specifically transformer-based models—to enhance classification performance on the same datasets utilized in the previous study.}
\ctufitkeywordsCZE{transformery, BERT, DistillBERT, RoBERTa, GPT falešné zprávy, hluboké učení, zpracování přirozeného jazyka, binární klasifikace}
\ctufitkeywordsENG{transformers, BERT, DistillBERT, RoBERTa, GPT, fake news, deep learning, natural language processing, binary classification}